% Part: incompleteness
% Chapter: incompleteness-provability
% Section: first-incompleteness-thm

\documentclass[../../../include/open-logic-section]{subfiles}

\begin{document}

\olfileid{inc}{inp}{1in}

\olsection{د ناتکميلۍ لومړۍ قضيه}

اوس د ګوډل د ناتکميلۍ د لومړۍ قضيې اصلي ثبوت بيانولے شو۔ فرض کړئ
$\Th{T}$ په هغې ژبه کښې هره په محاسبوي ډول بديهي‌کړې تيوري ده چې
د حساب ژبه غځوي، او~$\Th{T}$ د~$\Th{Q}$ بديهي اصول رانغاړي۔ دا په
ځانګړي ډول دا معنا لري چې~$\Th{T}$ محاسبه کېدونکې تابعې او اړيکې
تمثيلوي۔

موږ استدلال کړے دے چې که فورمولونه او ثبوتونه په مناسب ډول د عددونو
په توګه کوډ شي، نو اړيکه~$\Prf[\Th{T}](x,y)$ محاسبه کېدونکې ده۔
$\Prf[\Th{T}](x,y)$ هله او يوازې هله برقرارېږي چې~$x$ په~$\Th{T}$
کښې د هغه !!{formula} د !!a{derivation} ګوډل شمېره وي چې ګوډل شمېره
ئې~$y$ ده۔ په حقيقت کښې ګوډل د خپلې ټاکلې تيورۍ دپاره وښودله چې
دا اړيکه ابتدايي بازګشتي ده؛ د دې دپاره ئې په خپله مقاله کښې د
۴۵ تابعو او اړيکو لړ وکاراوه۔ پينځه څلوېښتمه اړيکه، $x B y$، د
ګوډل د~$\Th{T}$ د هماغې ټاکنې دپاره همدا~$\Prf[\Th{T}](x,y)$ ده۔
ياد ولرئ چې ګوډل په خپله مقاله کښې «بازګشتي» وايي، خو موږ به اوس
ورته «ابتدايي بازګشتي» وايو۔

څرنګه چې~$\Prf[\Th{T}](x,y)$ محاسبه کېدونکې ده، په~$\Th{T}$ کښې
تمثيلېدے شي۔ د هغې تمثيلوونکي فورمول دپاره~$\OPrf[\Th{T}](x,y)$
کاروو۔ $\OProv[\Th{T}](y)$ دې فورمول
$\lexists[x][\OPrf[\Th{T}](x,y)]$ وي۔ دا د ګوډل په لړ کښې
شپږڅلوېښتمه اړيکه، $\fn{Bew}(y)$، بيانوي۔ لکه څنګه چې ګوډل يادونه
کوي، دا يوازينۍ اړيکه ده چې «بازګشتي نۀ شو بللے»۔ ښايي د هغه مطلب
دا و: له تعريفه ئې محاسبه کېدنه نۀ څرګندېږي؛ او وروستيو پرمختګونو
په حقيقت کښې وښودله چې محاسبه کېدونکې هم نۀ ده۔

فرض کړئ $\Th{T}$ داسې !!{axiomatizable} تيوري ده چې~$\Th{Q}$
رانغاړي۔ نو~$\Prf[\Th{T}](x, y)$ د پرېکړې وړ ده، او ځکه په~$\Th{Q}$
کښې د يوه !!a{formula}~$\OPrf[\Th{T}](x, y)$ په وسيله تمثيلېږي۔
$\OProv[\Th{T}](y)$ دې هغه فورمول وي چې پورته مو بيان کړ۔ د ثابت
ټکي د لمې له مخې داسې فورمول~$!G_\Th{T}$ شته چې~$\Th{Q}$، او ځکه
$\Th{T}$ هم، دا !!{derive}s کوي:
\begin{equation}
\ollabel{eqn:qpf}
!G_\Th{T} \liff \lnot \OProv[\Th{T}](\gn{!G_\Th{T}}).
\end{equation}
پام وکړئ چې~$!G_\Th{T}$ په اصل کښې وايي: «$!G_\Th{T}$ په~$\Th{T}$
کښې !!{derivable} نۀ دے»۔

\begin{lem}\ollabel{lem:cons-G-unprov}
که~$\Th{T}$ سازګاره او !!{axiomatizable} تيوري وي چې~$\Th{Q}$
غځوي، نو $\Th{T} \Proves/ !G_\Th{T}$۔
\end{lem}

\begin{proof}
فرض کړئ $\Th{T}$، $!G_\Th{T}$ !!{derive}s کوي۔ نو د هغه
!!a{derivation} \emph{شته} دے، او په دې توګه د کوم عدد~$m$ دپاره
اړيکه~$\Prf[\Th{T}](m,
\Gn{!G_\Th{T}})$ برقرارېږي۔ خو بيا~$\Th{Q}$ جمله
$\OPrf[\Th{T}](\num m, \gn{!G_\Th{T}})$ !!{derive}s کوي۔ نو~$\Th{Q}$
$\lexists[x][\OPrf[\Th{T}](x,\gn{!G_\Th{T}})]$ هم !!{derive}s کوي،
چې د تعريف له مخې~$\OProv[\Th{T}](\gn{!G_\Th{T}})$ دے۔ د
\olref{eqn:qpf} له مخې~$\Th{Q}$، $\lnot !G_\Th{T}$ !!{derive}s کوي؛
او ځکه چې~$\Th{T}$، $\Th{Q}$ غځوي، $\Th{T}$ ئې هم ثابتوي۔
موږ وښودله چې که~$\Th{T}$، $!G_\Th{T}$ !!{derive}s کړي، نو د هغه
نفي~$\lnot !G_\Th{T}$ هم !!{derive}s کوي؛ په دې صورت کښې به تيوري
ناسازګاره وي۔
\end{proof}

\begin{defn}
\ollabel{thm:oconsis-q}
تيوري~$\Th{T}$ \emph{$\omega$-سازګاره} ده که لاندې شرط پوره کړي:
که~$\lexists[x][!A(x)]$ هره جمله وي او~$\Th{T}$ د
$\lnot !A(\num 0)$، $\lnot !A(\num 1)$، $\lnot !A(\num 2)$،
\dots هره يوه !!{derive}s کړي، نو~$\Th{T}$،
$\lexists[x][!A(x)]$ نۀ ثابتوي۔
\end{defn}

پام وکړئ چې هره $\omega$-سازګاره تيوري سازګاره هم ده۔ دا نېغ له
دې خبرې راځي چې که~$\Th{T}$ ناسازګاره وي، نو د هر~$!A$ دپاره
$\Th{T} \Proves !A$۔ په ځانګړي ډول که~$\Th{T}$ ناسازګاره وي، د
هر~$n$ دپاره $\lnot !A(\num n)$ !!{derive}s کوي او ورسره
$\lexists[x][!A(x)]$ هم !!{derive}s کوي۔ نو ناسازګاره~$\Th{T}$
$\omega$-ناسازګاره ده۔ د متقابل عکس له مخې، که~$\Th{T}$
$\omega$-سازګاره وي، نو بايد سازګاره وي۔

\begin{lem}\ollabel{lem:omega-cons-G-unref}
که~$\Th{T}$، $\omega$-سازګاره او !!{axiomatizable} تيوري وي چې
$\Th{Q}$ غځوي، نو $\Th{T} \Proves/ \lnot !G_\Th{T}$۔
\end{lem}

\begin{proof}
ښيو چې که~$\Th{T}$، $\lnot !G_\Th{T}$ !!{derive}s کړي، نو
$\omega$-ناسازګاره ده۔ فرض کړئ~$\Th{T}$، $\lnot !G_\Th{T}$
!!{derive}s کوي۔ که~$\Th{T}$ ناسازګاره وي، نو $\omega$-ناسازګاره
ده، او ثبوت بشپړ دے۔ که نۀ، نو~$\Th{T}$ سازګاره ده او د
\olref{lem:cons-G-unprov} له مخې~$!G_\Th{T}$ نۀ !!{derive} کوي۔
څرنګه چې په~$\Th{T}$ کښې د~$!G_\Th{T}$ هېڅ !!{derivation} نشته،
نو~$\Th{Q}$ دا !!{derive}s کوي:
\[
\lnot \OPrf[\Th{T}](\num 0, \gn{!G_\Th{T}}), \lnot \OPrf[\Th{T}](\num 1,
\gn{!G_\Th{T}}), \lnot \OPrf[\Th{T}](\num 2, \gn{!G_\Th{T}}), \dots
\]
او~$\Th{T}$ هم همداسې کوي۔ له بلې خوا د~\olref{eqn:qpf} له مخې
$\lnot !G_\Th{T}$ له
$\lexists[x][\OPrf[\Th{T}](x,\gn{!G_\Th{T}})]$ سره معادل دے۔
نو~$\Th{T}$، $\omega$-ناسازګاره ده۔
\end{proof}

\begin{prob}
  هره $\omega$-سازګاره تيوري سازګاره ده۔ وښيئ چې برعکس نۀ ده؛
  يعنې داسې تيورۍ شته چې سازګارې خو $\omega$-ناسازګارې دي۔ دا په
  ښودلو ثابت کړئ چې $\Th{Q} \cup
  \{\lnot !G_\Th{Q}\}$ سازګاره خو $\omega$-ناسازګاره ده۔
\end{prob}

\begin{thm}
\ollabel{thm:first-incompleteness} فرض کړئ $\Th{T}$ هره
$\omega$-سازګاره، !!{axiomatizable} تيوري ده چې~$\Th{Q}$ غځوي۔ نو
$\Th{T}$ بشپړه نۀ ده۔
\end{thm}

\begin{proof}
  که~$\Th{T}$، $\omega$-سازګاره وي، نو سازګاره ده؛ ځکه د
  \olref{lem:cons-G-unprov} له مخې $\Th{T}
  \Proves/ !G_\Th{T}$۔ د~\olref{lem:omega-cons-G-unref} له مخې
  $\Th{T} \Proves/ \lnot !G_\Th{T}$۔ دا معنا لري چې~$\Th{T}$
  ناتکميله ده، ځکه نۀ~$!G_\Th{T}$ او نۀ~$\lnot !G_\Th{T}$
  !!{derive}s کوي۔
\end{proof}

\end{document}
